\subsection{问题四：五次补充实验设计}

\subsubsection{由前三问证据缺口确定实验职责}

前三问留下的关键问题不是“模型能否再次得到相同答案”，而是现有数据尚不能回答：
关键高收率区间内部怎样变化、当前最高观测能否复现、领先组合在新增同温条件下是否
仍占优势，以及严格低温边界附近是否存在更高实测条件。下面依次由前三问的证据边界
确定新增实验的职责。

\paragraph{高收率区间的分辨率缺口}

第一问表明，各组合只有5--7个温度点，直线只能概括已测范围内的平均变化，不能确定
局部转折位置。第三问中，A3在350、400和450~℃的实测收率依次为18.0333\%、
44.7281\%和43.1184\%，相邻温度均相差50~℃。因此，现有数据只能把400~℃识别为
最高观测点，不能说明其两侧怎样变化，也不能判断区间内部是否存在更高实测条件。
新增实验需要在350--400~℃和400--450~℃内各补一个温度点。

\paragraph{最高观测的重复性缺口}

附件1的114个组合—温度条件均无独立重复。A3在400~℃的收率虽为当前最高观测值，
却只比其450~℃观测高1.6096个百分点。前三问能够复算并比较这些记录，但无法从单次
观测估计实验波动，也不能判断这一差距能否在相同条件下复现。因此，需要独立重复
A3在400~℃的实验，并直接报告两次观测差。

\paragraph{组合比较的同温证据缺口}

第二问发现不同组合的升温幅度并不一致，因而不能把A3自身区间加密的结果直接解释为
它仍优于其他组合。A3与主要竞争组合A4在350--400~℃之间缺少新增同温观测；若只补测
A3，就无法区分局部优势来自组合差异还是所选温度。为此，需要在同一新增温度下同时
测量A3和A4。

\paragraph{严格低温边界的覆盖缺口}

第三问在 \(T<350~^\circ\mathrm{C}\) 的严格低温情景下推荐A2在325~℃的实测条件，
而A2的下一个观测温度350~℃不满足严格不等式。现有数据因而不能回答325--350~℃之间
是否存在收率更高且仍满足约束的条件。新增实验需要在该开区间内选择一个设备可执行
的温度点。

设新增实验 \(e\) 的催化剂组合和反应温度分别为 \(c_e\) 和 \(T_e\)，新测乙醇
转化率、C4 烯烃选择性和收率分别为 \(X_e,S_e,Y_e\)。每次实验仍按
式\eqref{eq:c4-yield}的同一口径，由 \(X_e\) 与 \(S_e\) 计算 \(Y_e\)。

五次实验均沿用附件1中对应组合的完整配方和装料方式，只改变或重复反应温度，其他
可控条件保持可比。表\ref{tab:q4-design}把上述四类证据缺口落实为五次具体实验；
一次实验可以同时提供区间细化与同温比较信息。

\begin{table}[H]
\centering
\caption{五次补充实验的条件与职责}
\label{tab:q4-design}
\begin{tabular}{clclp{7.0cm}}
\toprule
实验 & 组合 & 温度（℃） & 类型 & 主要职责 \\
\midrule
E1 & A3 & 375   & 新条件 & 细分350--400~℃区间，并与E5作同温比较 \\
E2 & A3 & 425   & 新条件 & 细分400--450~℃区间，判断400~℃附近的高收率区域 \\
E3 & A2 & 337.5 & 新条件 & 细分严格低温情景下的325--350~℃区间 \\
E4 & A3 & 400   & 独立重复 & 检查当前最高观测条件的局部可复现性 \\
E5 & A4 & 375   & 新条件 & 与E1比较，检查A3对竞争组合的局部领先是否保持 \\
\bottomrule
\end{tabular}
\end{table}

A2、A3和A4均采用装料方式I。A2使用200~mg、2~wt\% Co/SiO\(_2\)，200~mg HAP和
1.68~mL/min乙醇条件；A3使用200~mg、1~wt\% Co/SiO\(_2\)，200~mg HAP和
0.9~mL/min乙醇条件；A4使用200~mg、0.5~wt\% Co/SiO\(_2\)，200~mg HAP和
1.68~mL/min乙醇条件。
这些条件均来自附件1，不创造或重组新配方。

\subsubsection{新增观测对既有结论的更新规则}

E1和E2分别与A3已有的350、400、450~℃观测比较。若新增条件取得更高实测收率，则
一般情景推荐相应更新；若没有超过原推荐，它们仍可缩小400~℃两侧的未知区间。
E3严格低于350~℃，若其收率超过A2在325~℃的17.2643\%，则低温推荐更新为
A2在337.5~℃；否则保留原推荐。

设A3在400~℃的原收率为 \(Y_0\)，E4的新测收率为 \(Y_4\)。该条件后续以

\begin{equation}
\overline{Y}_{\mathrm{A3},400}
=\frac{Y_0+Y_4}{2}
\tag{21}
\label{eq:q4-replicate-mean}
\end{equation}

作描述性排序，并同时报告 \(Y_4-Y_0\) 及其绝对值。由于其他条件仍多为单次观测，
式\eqref{eq:q4-replicate-mean}不能用于显著性推断或建立全域误差分布。

E1与E5在375~℃直接比较转化率、选择性和收率。若A4高于A3，只说明A3在这一局部
温度并非始终领先；若A3较高，也不能推广为所有温度下A3均优于A4。这两个新增点不
构成21种组合的共同温度表，因此只补充第二问，不替代其四个共同温度的主体分析。

五次实验完成后，一般情景在全部已有及新增实测条件中重新枚举；严格低温情景只枚举
\(T<350~^\circ\mathrm{C}\) 的实测条件。A3在400~℃按两次观测均值表示，插值预测
不进入实测排序。

\subsubsection{职责充分性与执行边界}

五次实验分别承担一般情景低侧加密、高侧加密、严格低温加密、关键高值复现和竞争组合
同温比较。逐项删除检查表明，删去E1会同时失去350--400~℃加密和同温竞争，删去
E2、E3、E4或E5则分别失去高侧加密、低温加密、局部可复现性检查或竞争组合比较。
因此该方案
在5次预算下职责紧凑，但这不表示所选温度在数学上唯一最优。

若实验结果能够在安排后续实验前获得，可先执行E1--E3，再把第四次用于重复更新后的
A3最高实测条件；第五次在低温推荐发生改变时重复更新后的A2低温最高条件，否则执行
A4在375~℃的实验。该序贯方案只在确有结果反馈权限时适用，不能与表
\ref{tab:q4-design}合并成7次实验。

337.5~℃来自325和350~℃的等距中点。正式执行前须确认设备能够稳定设置该温度；
若设备不支持，必须按真实控温精度重新选择区间内部温度，不能静默舍入后继续声称等距
二分。实验还应记录实际温度、批次和执行顺序，以免把设备或时间漂移误解释为组合差异。

受5次预算限制，该方案不能补齐21种组合在325、375、425或450~℃的完整交叉表，
不能为全部114个原条件增加重复，也不能估计长期失活、全域随机误差或工业放大效应。
这些未覆盖缺口继续构成结论边界；本方案只优先补充最可能改变当前推荐的局部证据。

综上，五次实验由前三问的证据缺口直接推出：第一问确定局部加密和重复位置，第二问
确定竞争组合的同温比较，第三问确定优先细化的两种推荐区域。该设计用于补足决策信息，
不把未来观测预写为结论，也不宣称五次实验足以建立完整响应面或用于证明前三问正确。
